\documentclass[11pt]{article}
\usepackage[margin=1in]{geometry}
\usepackage[T1]{fontenc}
\usepackage{enumitem}
\setlength{\parindent}{0pt}
\setlength{\parskip}{6pt}

\title{LibraryCard --- Friday Questions}
\author{}
\date{}

\begin{document}
\maketitle

\begin{enumerate}[leftmargin=2em, itemsep=1em]

\item \textbf{Why must the card number get its value in the member
initialization list?}

\texttt{number} is declared as \texttt{const int}, and a \texttt{const} data
member cannot be assigned after the object exists --- every assignment to it
in the constructor body is a compile error. The member-initializer list is the
one place where the member is \emph{born with} its value instead of being
default-constructed and then overwritten, so \texttt{number(nextNumber++)}
must appear there. A side effect is that the compiler, not my memory, is what
guarantees the number never changes afterwards.

\item \textbf{Why does \texttt{printReport(const LibraryCard\& card)} force you
to mark some member functions const?}

The \texttt{const} on the reference promise the function will not modify the
card it receives, so inside \texttt{printReport} the parameter is treated as a
const object. C++ only lets a const object call \emph{const member functions},
because only those promise (and the compiler enforces) that they do not change
the object's state. If \texttt{getHolder()} or \texttt{print()} were not
marked \texttt{const}, the calls \texttt{card.getHolder()} and
\texttt{card.print()} would fail to compile, since those functions could in
principle write to the card. So marking every read-only member function
\texttt{const} is not optional decoration --- it is what makes the card usable
through a read-only reference at all.

\item \textbf{At the end of main, \texttt{lobby[1]}, \texttt{lobby[0]} and
\texttt{ana} all close. In what order, and why?}

The order is \texttt{lobby[1]} (card 5006) first, then \texttt{lobby[0]}
(card 5005), and \texttt{ana} (card 5001) last of those three. Automatic
(local) objects are destroyed in \emph{reverse order of construction} when
their block ends. \texttt{ana} was constructed first of all the locals, so it
is destroyed last; the array \texttt{lobby} was constructed last, so it is
destroyed first, and its elements go in reverse index order
(\texttt{lobby[1]} before \texttt{lobby[0]}). Between them, \texttt{x} (5003)
and \texttt{ben} (5002) close in the same reverse pattern, which is exactly
why the output ends with \texttt{5006, 5005, 5003, 5002, 5001}.

\item \textbf{How can \texttt{LibraryCard::activeCards()} be called when no
card exists?}

Because \texttt{activeCards()} is a \texttt{static} member function: it
belongs to the class itself rather than to any one card object, so it does not
need an instance to be called --- \texttt{LibraryCard::activeCards()} invokes
it directly on the class. What it reads is also static: the \texttt{count}
data member lives once for the whole program, not inside each card, so it
exists even at the start of \texttt{main} when zero cards have been created.
Constructors increment \texttt{count} and the destructor decrements it, so the
static function always reports the current number of live cards.

\item \textbf{Reflection.}

This exercise tied the whole week's lecture material together in one small
class, and building it forced me to understand each feature as a tool rather
than syntax. The biggest ``aha'' was \texttt{const}: I initially thought of it
as labeling data that never changes, but the \texttt{printReport} question
showed me const is really a promise about member functions --- a contract
between the class and its callers. Delegating constructors and the private
static helpers were the other win; seeing all three construction paths funnel
through one body made the ``write each rule once'' rule feel natural instead
of dogmatic. The deleted copy operations and the destructor output also made
object lifetime concrete for the first time --- I could \emph{see} the reverse
destruction order in the program's output instead of just being told it
exists.

Discussing the homework with the other members of my home group helped too:
walking through why \texttt{temp}'s destructor fires early, and comparing our
different answers for the must-not-compile checks, surfaced details I had
glossed over (for example, why the string constructor has to be
\texttt{explicit} rather than just non-copyable). My one point of feedback for
the week is that the assignment's habit of stating behaviors without naming
the C++ feature was frustrating while I was stuck but probably the most
valuable part afterwards --- being forced to pick the right mechanism myself
is what made the features stick. I would not change that, but a short
checklist of ``which week-6 feature could produce this behavior'' after the
due date might help students verify they found the intended solution.

\end{enumerate}

\end{document}
